% preamble.tex -- shared by every chapter of "Next.js 성능: 세 글로 본 우리 앱 점검".
% Build (clean first, latex.md §12): latexmk -C && latexmk -xelatex main.tex
%
% Checklist item (one audit row), added for this book:
%   \begin{checkitem}{\ckdone|\ckopen|\ckna}{A1}{title}{STATUS · 영향 M}
%     \item[What] ... \item[Why] ... \item[How] ... \item[When] ...
%   \end{checkitem}
%   The four header arguments are macro arguments: use \texttt, never \code, in them.
% Chapter labels: ch:overview ch:sources ch:dynamic ch:actions ch:suspense
%   ch:cross ch:plan ch:results ch:lessons ch:checklist
%
% ===== MACROS (inherited from the traffic-security book) ====================
% Text macros (semantic names, never change them in the body):
%   \term{x}        a defined or named concept (bold). First use: gloss with \eng.
%   \eng{x}         English gloss in parentheses, no space: 캐시 키\eng{cache key}
%   \code{x}        command, path, header, identifier. Breaks at / _ . = etc.
%                   Do NOT use inside \caption, \section or another macro argument;
%                   use \texttt{} there.
%   \decision{N}    prints "결정 #N" (an owner decision row of decisions.plan.md)
%   \commit{hash}   a git commit hash in \texttt
%   \cfs{HIT}       a cf-cache-status value (HIT, MISS, DYNAMIC, BYPASS, ...)
%   \ctlmark{B}     a small red circle with the control letter (A..N), for text
%   \todo{x}        a gap the sources do not fill (red, boxed). Never fill it by guess.
% Environments:
%   callout[title]  grey box for a warning or a key sentence (latex.md §6)
%   inbrief         "In brief." with \item[Why] \item[When] \item[Where] \item[How]
% Column types (tabularx / tabular): Y = ragged-right X column,
%   L{w} = ragged-right paragraph column of width w.
% Colours: cfcol (Cloudflare, orange), vmcol (VM, blue), dbcol (Supabase, green),
%   badcol (failure, red), goodcol (pass, green), mutecol (baseline, grey).
% TikZ styles (on top of figkit's flowbox/flowalt/flownew/flowarrow):
%   cfbox (Cloudflare box), vmbox (VM box), dbbox (Supabase box),
%   zone (dashed container), ctl (red circle control marker),
%   lbl (small arrow label), badbox (failure box), goodbox (pass box).
% Chart style: \pgfplotsset{barchart} gives a bar chart that follows figures.md.
% Caption words are "Figure" and "Table" (latex.md §10). Labels: fig: tab: sec: ch:.
% ==========================================================================

\usepackage[a4paper,margin=25mm]{geometry}
\usepackage{fontspec}
\usepackage{xeCJK}
\setCJKmainfont{Apple SD Gothic Neo}   % latex.md §4
\xeCJKsetup{CJKspace=true}             % latex.md §2, required

\usepackage{amsmath}
\usepackage{booktabs}
\usepackage{tabularx}
\usepackage{longtable}
\usepackage{array}
\usepackage{etoolbox}
\usepackage{enumitem}
\usepackage{listings}
\usepackage[obeyspaces,spaces]{url}
\usepackage{microtype}
\usepackage{pifont}                    % check mark for checkitem
\usepackage{figkit}                    % float [H], caption, tikz, flowbox styles
\usetikzlibrary{fit,backgrounds,calc,shapes.geometric}
\usepackage{pgfplots}
\pgfplotsset{compat=1.17}
\usepackage[hidelinks]{hyperref}
\usepackage{cleveref}

\setlength{\emergencystretch}{3em}
\setlength{\parindent}{0pt}
\setlength{\parskip}{0.5em}
\setcounter{secnumdepth}{3}            % latex.md §15
\setcounter{tocdepth}{3}

% --- caption words: always English (latex.md §10) -----------------------------
\renewcommand{\figurename}{Figure}
\renewcommand{\tablename}{Table}
\renewcommand{\contentsname}{목차}
\renewcommand{\appendixname}{부록}
\captionsetup{font=small,labelfont=bf,labelsep=period}
\captionsetup[table]{position=above,skip=6pt}
\captionsetup[figure]{position=below,skip=6pt}
\crefname{figure}{Figure}{Figures}
\Crefname{figure}{Figure}{Figures}
\crefname{table}{Table}{Tables}
\Crefname{table}{Table}{Tables}
\crefname{section}{§}{§}
\Crefname{section}{§}{§}
\crefname{chapter}{Chapter}{Chapters}
\Crefname{chapter}{Chapter}{Chapters}

% --- figkit legend words in Korean --------------------------------------------
\renewcommand{\flowaltname}{기준선}
\renewcommand{\flownewname}{바뀐 단계}

% --- table helpers -------------------------------------------------------------
\newcolumntype{Y}{>{\raggedright\arraybackslash}X}
\newcolumntype{L}[1]{>{\raggedright\arraybackslash}p{#1}}
\renewcommand{\arraystretch}{1.15}

% --- colours -------------------------------------------------------------------
\definecolor{cfcol}{RGB}{230,126,34}
\definecolor{vmcol}{RGB}{41,98,180}
\definecolor{dbcol}{RGB}{39,140,90}
\definecolor{badcol}{RGB}{192,57,43}
\definecolor{goodcol}{RGB}{39,140,90}
\definecolor{mutecol}{RGB}{120,120,120}

% --- semantic text macros ------------------------------------------------------
\DeclareUrlCommand\codeurl{\urlstyle{tt}}
\newcommand{\code}[1]{\allowbreak\hspace{0.12em}\codeurl{#1}\hspace{0.12em}\allowbreak}
\newcommand{\ckdot}{\hspace{0.3em}\textperiodcentered\hspace{0.3em}}
\newcommand{\term}[1]{\textbf{#1}}
\newcommand{\eng}[1]{(#1)}
\newcommand{\decision}[1]{결정~\#{}#1}
\newcommand{\commit}[1]{\texttt{#1}}
\newcommand{\cfs}[1]{\texttt{#1}}
\newcommand{\todo}[1]{\textcolor{badcol}{\textbf{[TODO: #1]}}}

\tikzset{
  ctl/.style     = {circle, draw=badcol, fill=badcol!12, text=black, inner sep=1pt,
                    minimum size=1.1em, font=\scriptsize\bfseries},
  cfbox/.style   = {flowbox, draw=cfcol, fill=cfcol!10},
  vmbox/.style   = {flowbox, draw=vmcol, fill=vmcol!8},
  dbbox/.style   = {flowbox, draw=dbcol, fill=dbcol!10},
  badbox/.style  = {flowbox, draw=badcol, fill=badcol!8},
  goodbox/.style = {flowbox, draw=goodcol, fill=goodcol!10},
  zone/.style    = {draw=black!45, dashed, rounded corners=4pt, inner sep=6pt},
  lbl/.style     = {font=\scriptsize, inner sep=1pt, fill=white, align=center},
}
\newcommand{\ctlmark}[1]{\tikz[baseline=(c.base)]\node[ctl] (c) {#1};}

% --- chart style (figures.md: white background, light grid, legend above) -------
\pgfplotsset{
  barchart/.style = {
    ybar stacked, bar width=14pt, ymajorgrids=true, grid style={black!12},
    axis lines=left, tick label style={font=\footnotesize},
    label style={font=\footnotesize}, legend style={font=\footnotesize, draw=none},
    legend columns=-1, legend style={at={(0.5,1.04)}, anchor=south},
    enlarge x limits=0.15, ymin=0,
  },
}

% --- callout box without tcolorbox (latex.md §6) --------------------------------
\newsavebox{\calloutbox}
\newenvironment{callout}[1][]
  {\par\medskip\noindent
   \begin{lrbox}{\calloutbox}\begin{minipage}{0.92\linewidth}
   \def\calloutttl{#1}\ifx\calloutttl\empty\else\textbf{#1}\par\smallskip\fi}
  {\end{minipage}\end{lrbox}%
   \fcolorbox{black!30}{black!3}{\usebox{\calloutbox}}\par\medskip}

% --- In brief (book.md principle 7) ---------------------------------------------
\newenvironment{inbrief}
  {\paragraph{In brief.}\begin{description}[font=\bfseries]}
  {\end{description}}

% --- listings: never split across pages (latex.md §9) ----------------------------
\lstset{basicstyle=\ttfamily\scriptsize, breaklines=true, columns=fullflexible,
        keepspaces=true, frame=single, rulecolor=\color{black!30}}
\BeforeBeginEnvironment{lstlisting}{\par\medskip\noindent\begin{minipage}{\linewidth}}
\AfterEndEnvironment{lstlisting}{\end{minipage}\par\medskip}

% --- checklist item (audit row) -------------------------------------------------
% Glyphs are drawn, not typed: Apple SD Gothic Neo has no ballot-box glyph.
\newcommand{\ckbox}[1]{\fbox{\makebox[0.9em][c]{\strut #1}}}
\newcommand{\ckdone}{\ckbox{\textcolor{goodcol}{\ding{51}}}}
\newcommand{\ckopen}{\ckbox{\phantom{\ding{51}}}}
\newcommand{\ckna}{\ckbox{\scriptsize\textcolor{mutecol}{--}}}
\newenvironment{checkitem}[4]
  {\par\medskip\noindent\begin{minipage}{\linewidth}\raggedright
   #1\ \textbf{#2.\ #3}\par
   \hspace*{1.6em}{\small\textcolor{mutecol}{#4}}\end{minipage}\par\nopagebreak\smallskip
   \begin{description}[font=\bfseries\small, leftmargin=3.4em, labelwidth=3em,
                       itemsep=1pt, parsep=1pt, topsep=0pt, partopsep=0pt]\raggedright}
  {\end{description}\par\medskip}

% --- \code / \url may also break after a hyphen (long article URLs) ------------------
\expandafter\def\expandafter\UrlBreaks\expandafter{\UrlBreaks\do\-}
% keep "*" tight inside \code (url sets it as a math binary operator)
\expandafter\def\expandafter\UrlSpecials\expandafter{\UrlSpecials\do\*{\mbox{\ttfamily *}}}
